\documentclass[a4paper,twoside]{article}

\usepackage{morewrites}

\usepackage[svgnames,dvipsnames,x11names]{xcolor}
\usepackage{graphicx}
\usepackage[english]{babel}
\usepackage[colorlinks=true, linkcolor=NavyBlue, urlcolor=NavyBlue, citecolor=NavyBlue]{hyperref}
\usepackage[a4paper]{geometry}
\usepackage{ts-fonts}
\usepackage[normalem]{ulem}
\usepackage{booktabs}
\usepackage{tabularx}
\usepackage{amsmath}
\usepackage{float}
\usepackage{seqsplit}
\newcolumntype{Y}{>{\centering\arraybackslash}X}
\newcolumntype{Z}{>{\raggedleft\arraybackslash}X}
\usepackage{ts-code}

\usepackage{lastpage}
\usepackage{refcount}
\makeatletter
\def\ps@plain{
  \let\@oddhead\@empty
  \let\@evenhead\@empty
  \def\@oddfoot{
    \hfil
    \ifnum\getpagerefnumber{LastPage}>1\relax
      \thepage
    \fi
    \hfil}
  \let\@evenfoot\@oddfoot
}
\makeatother

\title{Notes on TeXSmith Extensions}
\author{}\date{}

\begin{document}

\maketitle

\thispagestyle{plain}
TeXSmith ships with laser-focused extensions to close the \LaTeX{}-sized gaps that stock Markdown leaves behind. They cooperate with MkDocs and MkDocs Material, tagging the HTML with enough metadata for the \LaTeX{} renderer to finish the heavy lifting.

\begin{itemize}
\item{} Index entries double as Lunr tags in MkDocs and page references in LaTeX.
\item{} Glossary entries and acronyms keep terminology consistent.
\item{} Citations and cross-references wire figures, tables, equations, and bibliography entries together.
\item{} Raw LaTeX fences and inline snippets let you sprinkle precise TeX without polluting the HTML build.

\end{itemize}
\section{Syntax}\label{syntax}

The syntax for these helpers follows a few rules:

\begin{enumerate}
\item Easy to type in Markdown.
\item Friendly with standard Markdown parsers.
\item Collision-free with existing extensions.
\item Minimal visual noise in the raw text.

\end{enumerate}
\subsection{Syntax shorthand}\label{syntax-shorthand}

\begin{description}
\item[{ \texttt{@[]} Smart references }] Expands to the proper label for the referenced object. Figures become \texttt{Figure X}, tables render as \texttt{Table Y}, sections as \texttt{Section Z}, equations as \texttt{(N)}, theorems as \texttt{Theorem M}, and so on.
\item[{ \texttt{\^{}[]} Footnotes/bibliographic citations }] Resolves to either an inline footnote (if defined locally) or a bibliography citation.
\item[{ \texttt{\{index\}[...]} Index entries }] Inserts one or more index terms. They stay invisible in HTML yet show up as tags in Lunr and as \texttt{\textbackslash{}index\{\}} entries in LaTeX.
\item[{ \texttt{\{margin\}[...]\{side?\}} Margin notes }] Renders as \texttt{\textbackslash{}marginnote\{…\}} from the LaTeX \texttt{marginnote} package (auto-loaded on first use). The optional suffix \texttt{\{l\}} / \texttt{\{r\}} / \texttt{\{o\}} / \texttt{\{i\}} forces a margin; without it the note follows the document's default (right in \texttt{oneside}, outer in \texttt{twoside}).

\end{description}
Quick reference:

\begin{itemize}
\item{} \texttt{[](:)} Add content block
\item{} \texttt{@()} Cross-reference helper
\item{} \texttt{\^{}[]} / \texttt{\textbackslash{}cite\{\}} Bibliographic citation
\item{} \texttt{\{index\}[term]} / \texttt{\textbackslash{}index\{\}} Index entry
\item{} \texttt{\{margin\}[note]\{l|r|o|i\}?} / \texttt{\textbackslash{}marginnote\{\}} Margin note

\end{itemize}
\section{Other extensions}\label{other-extensions}

\begin{itemize}
\item{} Epigraphs: define \texttt{epigraph} blocks in front matter or mark blockquotes with the \texttt{epigraph} class to route them through the LaTeX epigraph macro.
\item{} Drop caps: Material's \texttt{lettering} syntax (\texttt{:[A](Natoly)}) produces LaTeX lettrine output automatically.
\item{} Wikipedia shortcodes keep working; they simply render as annotated links in both HTML and LaTeX.

\end{itemize}
\section{Index}\label{index}

Printed indexes convey intent with typography:

\begin{itemize}
\item{} Normal text: the topic is discussed.
\item{} Italic: quick mention only.
\item{} Bold: this section focuses on the topic.
\item{} Bold italic: primary topic plus ancillary references.
\item{} Nested entries: group related terms.

\end{itemize}
Use \texttt{\{index\}} plus multiple \texttt{[brackets]} to build entries. Append \texttt{\{b\}}, \texttt{\{i\}}, or \texttt{\{bi\}} to tweak the style, and specify \texttt{\{index:registry\}} when you want to file the entry under a custom registry (handy for multi-index books).

\begin{code}{md}{}{}
\begin{Verbatim}[commandchars=\\\{\},breaklines, breakanywhere, commandchars=\\\{\}]
Do you know the Gulliver\PYZsq{}s Travels tale about the egg dispute?
\PYZob{}index\PYZcb{}[endianness]\PYZob{}i\PYZcb{}

\PYZob{}index\PYZcb{}[endianness]\PYZob{}bi\PYZcb{}
\PYZob{}index\PYZcb{}[endianness]\PYZob{}b\PYZcb{}
\PYZob{}index\PYZcb{}[byte order][endianness]\PYZob{}i\PYZcb{}
\PYZob{}index:physics\PYZcb{}[relativity]\PYZob{}b\PYZcb{}
\end{Verbatim}
\end{code}
\section{Margin notes}\label{margin-notes}

Drop short asides into the margin with the \texttt{\{margin\}[…]\{side?\}} inline syntax.
An optional single-letter suffix selects the side:

\begin{center}
\begin{tabularx}{\linewidth}{>{\raggedright\arraybackslash}X>{\raggedright\arraybackslash}X}
\toprule
\textbf{Suffix} & \textbf{Semantics} \\

\midrule
(none) & document default --- right in \texttt{oneside}, outer in \texttt{twoside} \\
\texttt{\{l\}} & force the left margin (scoped \texttt{\textbackslash{}reversemarginpar}) \\
\texttt{\{r\}} & force the right margin \\
\texttt{\{o\}} & outer --- MVP alias of \texttt{\{r\}} \\
\texttt{\{i\}} & inner --- MVP alias of \texttt{\{l\}} \\

\bottomrule
\end{tabularx}
\end{center}

\begin{code}{md}{}{}
\begin{Verbatim}[commandchars=\\\{\},breaklines, breakanywhere, commandchars=\\\{\}]
Hooke\PYZsq{}s law\PYZob{}margin\PYZcb{}[linear only at small strain] holds
below the yield point, but non\PYZhy{}linear effects\PYZob{}margin\PYZcb{}[see
\PY{g+gs}{**Prandtl 1921**} for the classical derivation]\PYZob{}l\PYZcb{} dominate
above it.
\end{Verbatim}
\end{code}
The note compiles down to a \texttt{\textbackslash{}marginnote\{…\}} call from the \LaTeX{} \texttt{marginnote}
package, which TeXSmith's \texttt{ts-\allowbreak{}extra} fragment auto-loads when it spots the
command in the rendered output. A forced left side wraps the call in a
group that locally applies \texttt{\textbackslash{}reversemarginpar}, so subsequent notes go
back to the default placement.

Inline Markdown inside the note (\texttt{**bold**}, \texttt{*italic*}, \texttt{`code`},
\texttt{[links](…)}) is preserved and runs through the standard inline handlers
on the way to \LaTeX{}.

\subsection{Width}\label{width}

Margin notes never bleed past the page edge. \texttt{ts-\allowbreak{}extra} ships three
defensive layers alongside \texttt{\textbackslash{}usepackage\{marginnote\}}:

\begin{enumerate}
\item \textbf{Geometry-aware clamp.} An \texttt{\textbackslash{}AtBeginDocument} hook clamps
   \texttt{\textbackslash{}marginparwidth} to whatever horizontal space the document's geometry
   actually reserves for the margin --- the smaller of the recto outer
   margin and, in \texttt{twoside} documents, the verso outer margin as well ---
   minus \texttt{\textbackslash{}marginparsep} and a 6 mm safety buffer.
\item \textbf{Line-breaker tolerance.} \texttt{\textbackslash{}marginfont} includes \texttt{\textbackslash{}sloppy},
   \texttt{\textbackslash{}emergencystretch=1em} and \texttt{\textbackslash{}hyphenpenalty=50} so a long technical
   word, URL or German compound hyphenates (or stretches inter-word
   spacing) rather than poking past the margin box edge.
\item \textbf{Footnote-size body.} Notes default to \texttt{\textbackslash{}footnotesize} so multi-word
   notes still fit in narrow margins.

\end{enumerate}
That means a \texttt{geometry} block like

\begin{code}{yaml}{}{}
\begin{Verbatim}[commandchars=\\\{\},breaklines, breakanywhere, commandchars=\\\{\}]
\PY{n+nt}{press}\PY{p}{:}
\PY{+w}{  }\PY{n+nt}{geometry}\PY{p}{:}
\PY{+w}{    }\PY{n+nt}{left}\PY{p}{:}\PY{+w}{ }\PY{l+lScalar+lScalarPlain}{3cm}
\PY{+w}{    }\PY{n+nt}{right}\PY{p}{:}\PY{+w}{ }\PY{l+lScalar+lScalarPlain}{4cm}
\end{Verbatim}
\end{code}
automatically yields notes that fit on both recto and verso pages --- no
manual \texttt{marginparwidth=…} tweak needed.

\subsection{Font size}\label{font-size}

Because printed margins are narrow, margin notes default to \texttt{\textbackslash{}footnotesize}
so multi-word notes fit without overflowing. Override the default from your
own preamble (or a custom \texttt{ts-\allowbreak{}extra} snippet) if you want a different
treatment:

\begin{code}{latex}{}{}
\begin{Verbatim}[commandchars=\\\{\},breaklines, breakanywhere, commandchars=\\\{\}]
\PY{k}{\PYZbs{}renewcommand}\PY{k}{*}\PY{n+nb}{\PYZob{}}\PY{k}{\PYZbs{}marginfont}\PY{n+nb}{\PYZcb{}}\PY{n+nb}{\PYZob{}}\PY{k}{\PYZbs{}small}\PY{k}{\PYZbs{}itshape}\PY{n+nb}{\PYZcb{}}
\end{Verbatim}
\end{code}
\texttt{\textbackslash{}marginfont} is part of the \texttt{marginnote} package and is applied
automatically to every note.

\section{Citations}\label{citations}

Bibliographic references land in two ways:

\begin{enumerate}
\item Point TeXSmith at one or more \texttt{.bib} files and cite entries using \texttt{\^{}[]}.
\item Declare references directly in front matter via DOIs or inline metadata.

\end{enumerate}
\begin{code}{yaml}{}{}
\begin{Verbatim}[commandchars=\\\{\},breaklines, breakanywhere, commandchars=\\\{\}]
\PY{n+nn}{\PYZhy{}\PYZhy{}\PYZhy{}}
\PY{n+nt}{bibliography}\PY{p}{:}
\PY{+w}{  }\PY{c+c1}{\PYZsh{} Just DOI entry, TeXSmith will fetch the rest}
\PY{+w}{  }\PY{n+nt}{ein05}\PY{p}{:}\PY{+w}{ }\PY{l+lScalar+lScalarPlain}{https://doi.org/10.1002/andp.19053221004}
\PY{+w}{  }\PY{c+c1}{\PYZsh{} Manual entry}
\PY{+w}{  }\PY{n+nt}{KOFINAS2025}\PY{p}{:}
\PY{+w}{    }\PY{n+nt}{type}\PY{p}{:}\PY{+w}{ }\PY{l+lScalar+lScalarPlain}{article}
\PY{+w}{    }\PY{n+nt}{title}\PY{p}{:}\PY{+w}{ }\PY{p+pIndicator}{|}
\PY{+w}{        }\PY{n+no}{The impact of generative AI on academic integrity of authentic}
\PY{+w}{        }\PY{n+no}{assessments within a higher education context}
\PY{+w}{    }\PY{n+nt}{authors}\PY{p}{:}
\PY{+w}{      }\PY{p+pIndicator}{\PYZhy{}}\PY{+w}{ }\PY{n+nt}{name}\PY{p}{:}\PY{+w}{ }\PY{l+s}{\PYZdq{}}\PY{l+s}{Alexander}\PY{n+nv}{ }\PY{l+s}{K.}\PY{n+nv}{ }\PY{l+s}{Kofinas}\PY{l+s}{\PYZdq{}}
\PY{+w}{        }\PY{n+nt}{affiliation}\PY{p}{:}\PY{+w}{ }\PY{l+s}{\PYZdq{}}\PY{l+s}{University}\PY{n+nv}{ }\PY{l+s}{of}\PY{n+nv}{ }\PY{l+s}{Example}\PY{l+s}{\PYZdq{}}
\PY{+w}{      }\PY{p+pIndicator}{\PYZhy{}}\PY{+w}{ }\PY{l+s}{\PYZdq{}}\PY{l+s}{Crystal}\PY{n+nv}{ }\PY{l+s}{Han\PYZhy{}Huei}\PY{n+nv}{ }\PY{l+s}{Tsay}\PY{l+s}{\PYZdq{}}
\PY{+w}{      }\PY{p+pIndicator}{\PYZhy{}}\PY{+w}{ }\PY{l+s}{\PYZdq{}}\PY{l+s}{David}\PY{n+nv}{ }\PY{l+s}{Pike}\PY{l+s}{\PYZdq{}}
\PY{+w}{    }\PY{n+nt}{journal}\PY{p}{:}\PY{+w}{ }\PY{l+s}{\PYZdq{}}\PY{l+s}{British}\PY{n+nv}{ }\PY{l+s}{Journal}\PY{n+nv}{ }\PY{l+s}{of}\PY{n+nv}{ }\PY{l+s}{Educational}\PY{n+nv}{ }\PY{l+s}{Technology}\PY{l+s}{\PYZdq{}}
\PY{+w}{    }\PY{n+nt}{date}\PY{p}{:}\PY{+w}{ }\PY{l+lScalar+lScalarPlain}{2025\PYZhy{}03}
\PY{+w}{    }\PY{n+nt}{volume}\PY{p}{:}\PY{+w}{ }\PY{l+lScalar+lScalarPlain}{56}
\PY{+w}{    }\PY{n+nt}{number}\PY{p}{:}\PY{+w}{ }\PY{l+lScalar+lScalarPlain}{6}
\PY{+w}{    }\PY{n+nt}{pages}\PY{p}{:}\PY{+w}{ }\PY{l+s}{\PYZdq{}}\PY{l+s}{2522\PYZhy{}2549}\PY{l+s}{\PYZdq{}}
\PY{+w}{    }\PY{n+nt}{url}\PY{p}{:}\PY{+w}{ }\PY{l+lScalar+lScalarPlain}{https://doi.org/10.1111/bjet.13585}
\PY{n+nn}{\PYZhy{}\PYZhy{}\PYZhy{}}
\PY{l+lScalar+lScalarPlain}{We}\PY{l+lScalar+lScalarPlain}{ }\PY{l+lScalar+lScalarPlain}{know}\PY{l+lScalar+lScalarPlain}{ }\PY{l+lScalar+lScalarPlain}{that}\PY{l+lScalar+lScalarPlain}{ }\PY{l+lScalar+lScalarPlain}{time}\PY{l+lScalar+lScalarPlain}{ }\PY{l+lScalar+lScalarPlain}{is}\PY{l+lScalar+lScalarPlain}{ }\PY{l+lScalar+lScalarPlain}{relative}\PY{l+lScalar+lScalarPlain}{ }\PY{l+lScalar+lScalarPlain}{\PYZca{}[ein05]}\PY{l+lScalar+lScalarPlain}{ }\PY{l+lScalar+lScalarPlain}{and}\PY{l+lScalar+lScalarPlain}{ }\PY{l+lScalar+lScalarPlain}{recent}\PY{l+lScalar+lScalarPlain}{ }\PY{l+lScalar+lScalarPlain}{work}\PY{l+lScalar+lScalarPlain}{ }\PY{l+lScalar+lScalarPlain}{explores}
\PY{l+lScalar+lScalarPlain}{assessment}\PY{l+lScalar+lScalarPlain}{ }\PY{l+lScalar+lScalarPlain}{\PYZca{}[KOFINAS2025].}\PY{l+lScalar+lScalarPlain}{ }\PY{l+lScalar+lScalarPlain}{You}\PY{l+lScalar+lScalarPlain}{ }\PY{l+lScalar+lScalarPlain}{can}\PY{l+lScalar+lScalarPlain}{ }\PY{l+lScalar+lScalarPlain}{also}\PY{l+lScalar+lScalarPlain}{ }\PY{l+lScalar+lScalarPlain}{cite}\PY{l+lScalar+lScalarPlain}{ }\PY{l+lScalar+lScalarPlain}{multiple}
\PY{l+lScalar+lScalarPlain}{references}\PY{l+lScalar+lScalarPlain}{ }\PY{l+lScalar+lScalarPlain}{\PYZca{}[ein05,KOFINAS2025].}
\end{Verbatim}
\end{code}
Or with the CLI:

\begin{code}{sh}{}{}
\begin{Verbatim}[commandchars=\\\{\},breaklines, breakanywhere, commandchars=\\\{\}]
texsmith\PY{+w}{ }article.md\PY{+w}{ }article.bib
\end{Verbatim}
\end{code}
Or directly in Python:

\begin{code}{python}{}{}
\begin{Verbatim}[commandchars=\\\{\},breaklines, breakanywhere, commandchars=\\\{\}]
\PY{k+kn}{from}\PY{+w}{ }\PY{n+nn}{pathlib}\PY{+w}{ }\PY{k+kn}{import} \PY{n}{Path}

\PY{k+kn}{from}\PY{+w}{ }\PY{n+nn}{texsmith}\PY{+w}{ }\PY{k+kn}{import} \PY{n}{ConversionRequest}\PY{p}{,} \PY{n}{ConversionService}

\PY{n}{service} \PY{o}{=} \PY{n}{ConversionService}\PY{p}{(}\PY{p}{)}
\PY{n}{request} \PY{o}{=} \PY{n}{ConversionRequest}\PY{p}{(}
    \PY{n}{documents}\PY{o}{=}\PY{p}{[}\PY{n}{Path}\PY{p}{(}\PY{l+s+s2}{\PYZdq{}}\PY{l+s+s2}{article.md}\PY{l+s+s2}{\PYZdq{}}\PY{p}{)}\PY{p}{]}\PY{p}{,}
    \PY{n}{bibliography\PYZus{}files}\PY{o}{=}\PY{p}{[}\PY{n}{Path}\PY{p}{(}\PY{l+s+s2}{\PYZdq{}}\PY{l+s+s2}{article.bib}\PY{l+s+s2}{\PYZdq{}}\PY{p}{)}\PY{p}{]}
\PY{p}{)}
\PY{n}{response} \PY{o}{=} \PY{n}{service}\PY{o}{.}\PY{n}{execute}\PY{p}{(}\PY{n}{request}\PY{p}{)}
\PY{n}{tex\PYZus{}path} \PY{o}{=} \PY{n}{response}\PY{o}{.}\PY{n}{render\PYZus{}result}\PY{o}{.}\PY{n}{main\PYZus{}tex\PYZus{}path}
\PY{n+nb}{print}\PY{p}{(}\PY{l+s+sa}{f}\PY{l+s+s2}{\PYZdq{}}\PY{l+s+s2}{LaTeX written to: }\PY{l+s+si}{\PYZob{}}\PY{n}{tex\PYZus{}path}\PY{l+s+si}{\PYZcb{}}\PY{l+s+s2}{\PYZdq{}}\PY{p}{)}
\end{Verbatim}
\end{code}
\section{Math}\label{math}

Inline math uses the standard Markdown-friendly delimiters: \texttt{\$...\$} or \texttt{\textbackslash{}(...\textbackslash{})} for inline spans, \texttt{\$\$...\$\$} or \texttt{\textbackslash{}[...\textbackslash{}]} for display mode. Most Markdown engines (including MkDocs) pass these through untouched.

Numbered equations are not part of Markdown itself, so TeXSmith handles the heavy lifting for you.

\begin{code}{md}{}{}
\begin{Verbatim}[commandchars=\\\{\},breaklines, breakanywhere, commandchars=\\\{\}]
\PYZob{}\PYZsh{}pythagoras\PYZcb{}
: \PYZdl{}\PYZdl{}a\PYZca{}2 + b\PYZca{}2 = c\PYZca{}2\PYZdl{}\PYZdl{}

From @[pythagoras], we know that...
\end{Verbatim}
\end{code}
The equation receives an automatic number and can be referenced inline.

\section{Theorems}\label{theorems}

Use admonitions to define theorems, lemmas, definitions, and friends.

\begin{code}{md}{}{}
\begin{Verbatim}[commandchars=\\\{\},breaklines, breakanywhere, commandchars=\\\{\}]
!!! theorem \PYZdq{}Pythagorean Theorem\PYZdq{} \PYZob{}\PYZsh{}thm:pythagoras\PYZcb{}
    This is a theorem about right triangles and can be summarized in the next
    equation
    \PYZdl{}\PYZdl{} x\PYZca{}2 + y\PYZca{}2 = z\PYZca{}2 \PYZdl{}\PYZdl{}
\end{Verbatim}
\end{code}
It will be rendered as:

\begin{code}{latex}{}{}
\begin{Verbatim}[commandchars=\\\{\},breaklines, breakanywhere, commandchars=\\\{\}]
\PY{k}{\PYZbs{}begin}\PY{n+nb}{\PYZob{}}theorem\PY{n+nb}{\PYZcb{}}[Pythagorean theorem]
\PY{k}{\PYZbs{}label}\PY{n+nb}{\PYZob{}}pythagorean\PY{n+nb}{\PYZcb{}}
This is a theorem about right triangles and can be summarized in the next
equation
\PY{l+s+sb}{\PYZbs{}[}\PY{n+nb}{ x}\PY{n+nb}{\PYZca{}}\PY{l+m}{2}\PY{n+nb}{ }\PY{o}{+}\PY{n+nb}{ y}\PY{n+nb}{\PYZca{}}\PY{l+m}{2}\PY{n+nb}{ }\PY{o}{=}\PY{n+nb}{ z}\PY{n+nb}{\PYZca{}}\PY{l+m}{2}\PY{n+nb}{ }\PY{l+s}{\PYZbs{}]}
\PY{k}{\PYZbs{}end}\PY{n+nb}{\PYZob{}}theorem\PY{n+nb}{\PYZcb{}}
\end{Verbatim}
\end{code}
TeXSmith automatically generates the following admonition types:

\begin{itemize}
\item{} Theorem (\texsmithEmoji{📐})
\item{} Corollary (\texsmithEmoji{🧾})
\item{} Lemma (\texsmithEmoji{📜})
\item{} Proof (\texsmithEmoji{🔍})

\end{itemize}
\section{Glossary}\label{glossary}

Specific terms live in the glossary; shorthand belongs in the acronym list. Keep them separate so TeXSmith can decide when to expand, hyperlink, or index each one.

\begin{description}
\item[{ Glossary entries }] Definitions for full terms, whether single words or multi-word concepts.
\item[{ Acronyms }] Shortened forms such as NASA or UNESCO, often with an expanded description.

\end{description}
Define both collections in front matter:

\begin{code}{yaml}{}{}
\begin{Verbatim}[commandchars=\\\{\},breaklines, breakanywhere, commandchars=\\\{\}]
\PY{n+nt}{acronyms}\PY{p}{:}
\PY{+w}{  }\PY{n+nt}{nasa}\PY{p}{:}
\PY{+w}{    }\PY{n+nt}{name}\PY{p}{:}\PY{+w}{ }\PY{l+lScalar+lScalarPlain}{NASA}
\PY{+w}{    }\PY{n+nt}{description}\PY{p}{:}\PY{+w}{ }\PY{l+lScalar+lScalarPlain}{National}\PY{l+lScalar+lScalarPlain}{ }\PY{l+lScalar+lScalarPlain}{Aeronautics}\PY{l+lScalar+lScalarPlain}{ }\PY{l+lScalar+lScalarPlain}{and}\PY{l+lScalar+lScalarPlain}{ }\PY{l+lScalar+lScalarPlain}{Space}\PY{l+lScalar+lScalarPlain}{ }\PY{l+lScalar+lScalarPlain}{Administration}
\PY{+w}{  }\PY{n+nt}{unesco}\PY{p}{:}
\PY{+w}{    }\PY{n+nt}{name}\PY{p}{:}\PY{+w}{ }\PY{l+lScalar+lScalarPlain}{UNESCO}
\PY{+w}{    }\PY{n+nt}{description}\PY{p}{:}\PY{+w}{ }\PY{l+lScalar+lScalarPlain}{United}\PY{l+lScalar+lScalarPlain}{ }\PY{l+lScalar+lScalarPlain}{Nations}\PY{l+lScalar+lScalarPlain}{ }\PY{l+lScalar+lScalarPlain}{Educational,}\PY{l+lScalar+lScalarPlain}{ }\PY{l+lScalar+lScalarPlain}{Scientific}\PY{l+lScalar+lScalarPlain}{ }\PY{l+lScalar+lScalarPlain}{and}\PY{l+lScalar+lScalarPlain}{ }\PY{l+lScalar+lScalarPlain}{Cultural}\PY{l+lScalar+lScalarPlain}{ }\PY{l+lScalar+lScalarPlain}{Organization}
\PY{n+nt}{glossary}\PY{p}{:}
\PY{+w}{  }\PY{n+nt}{solid}\PY{p}{:}
\PY{+w}{    }\PY{n+nt}{name}\PY{p}{:}\PY{+w}{ }\PY{l+lScalar+lScalarPlain}{S.O.L.I.D.}
\PY{+w}{    }\PY{n+nt}{description}\PY{p}{:}\PY{+w}{ }\PY{p+pIndicator}{|}
\PY{+w}{        }\PY{n+no}{Acronym for five design principles intended to make software designs}
\PY{+w}{        }\PY{n+no}{more understandable, flexible, and maintainable.}

\PY{+w}{        }\PY{n+no}{1. Single Responsibility Principle}
\PY{+w}{        }\PY{n+no}{2. Open/Closed Principle}
\PY{+w}{        }\PY{n+no}{3. Liskov Substitution Principle}
\PY{+w}{        }\PY{n+no}{4. Interface Segregation Principle}
\PY{+w}{        }\PY{n+no}{5. Dependency Inversion Principle}
\PY{+w}{  }\PY{n+nt}{liskov}\PY{p}{:}
\PY{+w}{    }\PY{n+nt}{name}\PY{p}{:}\PY{+w}{ }\PY{l+lScalar+lScalarPlain}{Liskov}\PY{l+lScalar+lScalarPlain}{ }\PY{l+lScalar+lScalarPlain}{Substitution}\PY{l+lScalar+lScalarPlain}{ }\PY{l+lScalar+lScalarPlain}{Principle}
\PY{+w}{    }\PY{n+nt}{description}\PY{p}{:}\PY{+w}{ }\PY{p+pIndicator}{|}
\PY{+w}{        }\PY{n+no}{The Liskov Substitution Principle (LSP) states that objects of a}
\PY{+w}{        }\PY{n+no}{superclass should be replaceable with objects of a subclass without}
\PY{+w}{        }\PY{n+no}{affecting the correctness of the program. In other words, if S is a}
\PY{+w}{        }\PY{n+no}{subtype of T, then objects of type T in a program may be replaced with}
\PY{+w}{        }\PY{n+no}{objects of type S without altering any of the desirable properties of}
\PY{+w}{        }\PY{n+no}{that program (e.g., correctness).}
\end{Verbatim}
\end{code}
Inside the Markdown document you can reference glossary entries with the \texttt{gls:} prefix:

\begin{code}{md}{}{}
\begin{Verbatim}[commandchars=\\\{\},breaklines, breakanywhere, commandchars=\\\{\}]
From the well\PYZhy{}known [](gls:solid) principles, the following class must
be [](gls:liskov) Substitution Principle\PYZhy{}compliant.
\end{Verbatim}
\end{code}
\subsection{Wikipedia}\label{wikipedia}

Many glossary-worthy entries live on Wikipedia, and TeXSmith can pull their summaries automatically.

\begin{code}{md}{}{}
\begin{Verbatim}[commandchars=\\\{\},breaklines, breakanywhere, commandchars=\\\{\}]
From the well\PYZhy{}known [\PY{n+nt}{SOLID}](\PY{n+na}{https://en.wikipedia.org/wiki/SOLID})
\end{Verbatim}
\end{code}
When enabled, TeXSmith converts Wikipedia links into glossary entries for the printed document.

\begin{code}{toml}{}{}
\begin{Verbatim}[commandchars=\\\{\},breaklines, breakanywhere, commandchars=\\\{\}]
\PY{k}{[}\PY{k}{texsmith}\PY{k}{.}\PY{k}{extensions}\PY{k}{]}
\PY{n}{wikipedia\PYZus{}glossary}\PY{+w}{ }\PY{o}{=}\PY{+w}{ }\PY{k+kc}{true}
\end{Verbatim}
\end{code}
\section{Caption}\label{caption}

\begin{description}
\item[{ TeXSmith style }] Here is a figure with a caption:

\begin{code}{md}{}{}
\begin{Verbatim}[commandchars=\\\{\},breaklines, breakanywhere, commandchars=\\\{\}]
A diagram with 25\PYZpc{} width
: ![\PY{n+nt}{Diagram example}](\PY{n+na}{../examples/diagrams.png})\PYZob{}width=25\PYZpc{}\PYZcb{}

Table Caption Avec une grosse famille de chats  \PYZob{}\PYZsh{}bigcats\PYZcb{}
: | Cat Name    | Age | Color      |
  | \PYZhy{}\PYZhy{}\PYZhy{}\PYZhy{}\PYZhy{}\PYZhy{}\PYZhy{}\PYZhy{}\PYZhy{}\PYZhy{}\PYZhy{} | \PYZhy{}\PYZhy{}\PYZhy{}:| \PYZhy{}\PYZhy{}\PYZhy{}\PYZhy{}\PYZhy{}\PYZhy{}\PYZhy{}\PYZhy{}\PYZhy{}\PYZhy{} |
  | Whiskers    |  2  | Tabby      |
  | Mittens     |  5  | Black      |
\end{Verbatim}
\end{code}
\item[{ Pymarkdown style }] Here is a figure with a caption:

\begin{code}{md}{}{}
\begin{Verbatim}[commandchars=\\\{\},breaklines, breakanywhere, commandchars=\\\{\}]
![\PY{n+nt}{Diagram example}](\PY{n+na}{../examples/diagrams.png})\PYZob{}width=25\PYZpc{}\PYZcb{}

/// figure\PYZhy{}caption
    attrs: \PYZob{}\PYZsh{}foobar\PYZcb{}
    Avec un chocolat violet qui sent la \PY{g+gs}{**vanille**}
///
\end{Verbatim}
\end{code}

\end{description}
\section{Formatting}\label{formatting}

Plain Markdown covers \textbf{bold}, \emph{italic}, and \texttt{inline code}. PyMdown extensions add \sout{strikethrough}, \ifXeTeX
  \providecommand{\texsmithHighlight}[1]{\textcolor{yellow}{#1}}
  \providecommand{\hl}[1]{\textcolor{yellow}{#1}}
\else
  \ifLuaTeX
    \RequirePackage{lua-ul}
    \providecommand{\texsmithHighlight}[1]{
      \begingroup
      \setulcolor{yellow}
      \ul{#1}
      \endgroup
    }
    \providecommand{\hl}[1]{\texsmithHighlight{#1}}
  \else
    \providecommand{\texsmithHighlight}[1]{
      \begingroup
      \setlength{\fboxsep}{0.8pt}
      \colorbox{yellow}{#1}
      \endgroup
    }
  \fi
\fi

\texsmithHighlight{highlighted text}, \uline{inserted text}.

Small capitals are missing from the spec, so TeXSmith repurposes the double-underscore syntax for that effect:

\begin{code}{markdown}{}{}
\begin{Verbatim}[commandchars=\\\{\},breaklines, breakanywhere, commandchars=\\\{\}]
\PY{g+gs}{\PYZus{}\PYZus{}Small Capitals\PYZus{}\PYZus{}}
\end{Verbatim}
\end{code}
which renders as:

\begin{code}{latex}{}{}
\begin{Verbatim}[commandchars=\\\{\},breaklines, breakanywhere, commandchars=\\\{\}]
\PY{k}{\PYZbs{}textsc}\PY{n+nb}{\PYZob{}}Small Capitals\PY{n+nb}{\PYZcb{}}
\end{Verbatim}
\end{code}
\section{Tables}\label{tables}

One of the limitations of Markdown is the lack of support for complex table features such as multi-row and multi-column cells, cell alignment, and captions.

When a table is too large to fit on the page, try:

\begin{itemize}
\item{} Slightly resizing the table to fit the available width.
\item{} Allowing cells to wrap across multiple lines.
\item{} Rotating the table to landscape orientation.

\end{itemize}

\end{document}
